\section{Discussion}\label{sec:discussion}

\subsection{Separator information and the rate}

The upper bounds share a construction: a common kernel converts local
truncated functionals into laws whose separator marginals agree exactly.
For the preordering, positivity and normalization hold directly. For the
ordinary module, products of globally SOS kernels give positive densities
but have a mass defect. We restore exact normalization at the level of
signed densities and use one common correction to make them positive
(\cref{thm:mod}). Private conditional means, recourse repair, and
label-dependent corrections act inside bags while retaining this scalar
separator consistency. Additional polynomial constraints require a further
global feasible repair.

The examples show that finite separator information can obstruct convergence
even when local positivity is exact. In the paired two-bag models of
\cref{thm:quad-sharp,thm:affine-sharp}, the local minima are $-h$ and $h$,
and the exact local-measure gap is $2E_{2r}(h)$. This is a lower bound for
the finite SDP gap, rather than an identity for every sparse relaxation.
For the quadratic example, the order-one SDP gap is $1/4$, while the
local-measure gap is $3-2\sqrt2$; at order two they agree
(\cref{prop:exact-orders}). A Lipschitz derivative gives an inverse-square
approximation upper bound, attained in order by our piecewise quadratic
example. The function $\abs{x}$ has approximation error of exact order
$1/r$. Smoothness alone does not determine an exact rate: smoother functions
and polynomials can be approximated faster.

An active recourse constraint can create a convex corner, but activation
alone is not sufficient. For example, minimizing $y^2$ over
$y\in[-1,1]$, $y\ge x$, and $y\ge0$, with $x\in[-1,1]$, gives the value
$x_+^2$, whose derivative is continuous despite the active-set change at
zero. In the inverse-order example, by contrast, the value $\abs{x}$ has a
corner and the projected optimal multiplier jumps. The conditions of
\cref{thm:reg-cheb,thm:reg-holder} control the projected multipliers and
thereby improve the recourse repair estimate; they are not an equivalence
between constraint activation and nonsmoothness.

For the sharp inverse-order instance, merging bags or splitting the shared
domain at zero gives a low-order exact full-preordering certificate
(\cref{rem:merge-split}). Adding the separator statistic $\abs{x}$ to the
exact local-measure model also removes its gap. These instance-specific
remedies suggest a question for sparse optimization: how should a method
choose between increasing polynomial degree, enlarging bags, and exchanging
additional separator information?

\subsection{Computational interpretation and scope}

The exponent of $r$ in the recourse hierarchy's block sizes depends only
on the number of shared variables; at fixed shared width, its SDP dimensions
grow polynomially in the private dimension. This is a useful structural distinction from a
joint total-degree hierarchy. The constants in the accuracy bounds still
depend on objective coefficients, private dimension, fiber conditioning,
and multiplier regularity. The ordinary box bound is uniform in bag count
and tree shape after normalization by $\Cf$; its total SDP size grows with
the number of bags. A full preordering has up to $2^w$ blocks per bag,
whereas the ordinary module has $w+1$ in a bag of width $w$.

Discretization is a relevant comparator. For box objectives and fixed
private domains, Chebyshev grids and tree dynamic programming give
inverse-square error bounds. With exact local QP values and general affine
complete recourse, the corresponding bound has order $1/r$
(\cref{cor:pre-grid,rem:grid-baseline}). A grid optimum $G_r$ with proved
error $\varepsilon_r$ brackets the true optimum as
$G_r-\varepsilon_r\le\fstar\le G_r$. Thus discretization can provide lower
bounds as well as feasible solutions. Our results control the particular
sparse SDP, every feasible moment point, and its algebraic certificate
cone. The convergence bounds do not establish a generic numerical runtime
or a practical advantage over grid optimization. Solving local QPs,
constructing SDP tables, numerical conditioning, and certified extraction
from inexact moment data require their own analysis.

The recourse theorems assume continuous private variables belonging to one
bag, pointwise private convexity on the shared box, and fixed polytopes or
affine complete recourse with a fixed private constraint matrix. The
regularity theorem's inverse-square case uses weighted Chebyshev
summability or coordinatewise Lipschitz projected multipliers; its
coordinatewise Hölder case has rate $r^{-(1+\beta)}$ for $0<\beta\le1$.
The constrained extension requires a global distance bound, whose exponent
may be smaller than those of separate local geometric bounds. Its projection argument
proves existence of a feasible repaired law, rather than an efficient
projection algorithm.

Certificate conclusions depend on the model. The finite box SDPs have
attained real dual certificates. The private-degree-two hierarchy has the
same moment and certificate supremum values, with certificates at every
strict level below the moment optimum (\cref{prop:rec-dual}); boundary
attainment is not established. The polynomial-constraint hierarchies also
have equality of moment and certificate supremum values, and the proved
rates yield real certificates with any further positive slack
(\cref{prop:constr-dual}). This does not require strict feasibility or imply
boundary attainment. Exact rational box
certificates need positive slack beyond the finite-order error, and their
encoding length is polynomial in the expanded SDP dimensions and rational
input length, including the slack
(\cref{sec:certificates}). This size statement does not make those dimensions
polynomial in the original polynomial input.

\subsection{Open problems}

\begin{enumerate}
\item \emph{Logarithms for the ordinary module.} The general sparse upper
  bound is $\Cf\,O(\log^3r/r^2)$, while the quadratic example gives a lower
  bound of order $r^{-2}$. The necessity of these logarithms is unresolved
  by the present bounds.
\item \emph{Exact values beyond order two.} In the quadratic example of
  \cref{sec:sharpness}, the sparse SDP and exact local-measure values agree
  at order two. We conjecture agreement for every $r\ge2$. Retained
  floating-point computations are consistent with this conjecture but do
  not prove it. A proof would need a certificate for the relevant best
  separator approximation with a controlled degree.
\item \emph{Width dependence.} The displayed rates hold at fixed width.
  The ordinary-module correction has substantial width dependence; the
  correct dependence of the accuracy order on $w$ remains to be determined.
\item \emph{Ordinary modules with private recourse.} The recourse results
  use the shared box preordering. It remains to combine ordinary-module
  exact-consistency correction with conditional private means and recourse
  repair at comparable rates.
\item \emph{Multivariate multiplier regularity.} The proved hypotheses are
  weighted Chebyshev summability or coordinatewise dependence. Whether
  multivariate Lipschitz regularity suffices, including piecewise-affine
  projected multipliers with oblique region boundaries, remains open here.
\item \emph{Leading constants.} Bernstein's approximation constant for
  $\abs{x}$ determines the leading local-measure gap in the inverse-order
  example. We do not determine whether the finite SDP gaps have that same
  leading constant.
\item \emph{Boundary certificates.} The order-unit argument gives real
  recourse certificates with arbitrary strict slack. It does not settle
  membership in the certificate cone at the optimal moment level.
\end{enumerate}

The sparse preordering inverse-square rate was stated in lecture
slides~\cite{magron2025slides-lorentz,magron2026slides-tenors} before this
work. Our contribution for that cone is the complete compatible-density
proof and its explicit finite-order consequences. The ordinary-module
transfer and the recourse rate classification are separate theorem-specific
contributions, with the prior-work boundaries stated in the relevant
sections.
